% Generated by tools/edn_latex.py from reports/edn.md.
% Do not edit: change reports/edn.md and run `make edn`.
\ednentry{
  id = {EDN-02},
  title = {The model family is gradient boosting with LightGBM},
  date = {2026-09-22},
  milestone = {M1: Project Inception},
  topic = {Modelling Approach},
  participants = {All team members},
  decision = {Use gradient boosting, with LightGBM as the main model on a basic and an extended feature set and CatBoost as a challenger for high-cardinality categoricals.
A median and a Ridge baseline stay in the experiment ladder.},
  rationale = {\emph{Alternatives considered:}
\begin{itemize}[leftmargin=*, topsep=2pt]
\item Linear or Ridge regression only: simple and interpretable, but it misses interactions such as make by age by mileage.
\item kNN, random forest, tabular neural networks, hierarchical Bayes, LLM zero-shot: slower or weaker at this data size, or no fit for a CPU-only deployment with a small image.
\end{itemize}
\par
\emph{Why:} The course grades MLOps practice over raw accuracy.
LightGBM trains in minutes on a CPU, takes categoricals and missing values natively, keeps artefacts small, and gives exact SHAP explanations.
The baselines show whether it earns its added complexity.},
  ai = {Information seeking, Alternative assessment},
  response = {Accepted},
  assessment = {AI laid out the alternatives and the CPU, size, and explainability constraints.
The team reviewed them and confirmed the choice in sprint planning, including CatBoost as a challenger only.},
  aievidence = {``Why gradient boosting'' in \texttt{docs/\allowbreak{}docs/\allowbreak{}project-\allowbreak{}brief.\allowbreak{}md}, section 4 (2026-09-22).},
  evidence = {\href{https://github.com/mlops-2627q1-mds-upc/MLOPS_Recommenditos/pull/13}{PR \#13}.},
}
